\corrigesChapitre{13}{Circuits : intensités et tensions}
\par\noindent\begin{minipage}{\linewidth}
\begin{corrige}{1}{Le branchement}
\begin{questions}
\item En série dans la branche.
\item En dérivation aux bornes du dipôle.
\end{questions}
\pcFigureEx{figures/rev-cor-mesures.pdf}{L’ampèremètre est traversé par le courant de la branche ; le voltmètre est relié aux deux bornes du récepteur.}
\end{corrige}
\end{minipage}
\par\noindent\begin{minipage}{\linewidth}
\begin{corrige}{2}{Conversions}
\begin{questions}
\item $0{,}250\unite A$.
\item $32\unite{mA}$.
\end{questions}
\end{corrige}
\end{minipage}
\par\noindent\begin{minipage}{\linewidth}
\begin{corrige}{3}{En série}
\begin{questions}
\item $0{,}18\unite A$ en régime permanent.
\item Non ; la lampe reçoit et convertit de l’énergie.
\end{questions}
\pcFigureEx{figures/rev-cor-serie-018.pdf}{En régime permanent, la lampe ne consomme pas le courant ; l’intensité est la même avant et après.}
\end{corrige}
\end{minipage}
\par\noindent\begin{minipage}{\linewidth}
\begin{corrige}{4}{Deux nœuds}
\begin{questions}
\item En dérivation.
\item La tension, avec la même convention de mesure.
\end{questions}
\pcFigureEx{figures/rev-cor-parallele-12.pdf}{Deux branches reliées aux mêmes deux nœuds ont la même tension avec la même convention de mesure.}
\end{corrige}
\end{minipage}
\par\noindent\begin{minipage}{\linewidth}
\begin{corrige}{5}{Une somme}
\begin{questions}
\item $0{,}35\unite A$.
\item $350\unite{mA}$.
\end{questions}
\pcFigureEx{figures/rev-cor-noeud-035.pdf}{Somme des intensités entrantes égale à la somme des intensités sortantes.}
\end{corrige}
\end{minipage}
\par\noindent\begin{minipage}{\linewidth}
\begin{corrige}{6}{Une branche inconnue}
\begin{questions}
\item $I_2=0{,}80-0{,}250=0{,}550\unite A$.
\item On mélangerait A et mA.
\end{questions}
\par\smallskip\noindent\textbf{Pour comprendre.}\enspace La loi des nœuds donne $I_G=I_1+I_2$. Avec $I_1=0{,}250\unite A$, $I_2=0{,}80-0{,}250=0{,}550\unite A=550\unite{mA}$.\par
\pcFigureEx{figures/rev-cor-noeud-080.pdf}{250 mA = 0,250 A ; le courant de la branche inconnue est 0,550 A.}
\end{corrige}
\end{minipage}
\par\noindent\begin{minipage}{\linewidth}
\begin{corrige}{7}{Tensions en série}
\begin{questions}
\item $U_2=9-3{,}5=5{,}5\unite V$.
\item Aux bornes du second récepteur, en dérivation.
\end{questions}
\par\smallskip\noindent\textbf{Pour comprendre.}\enspace Les tensions des récepteurs en série s’ajoutent : $U_2=U_G-U_1=9-3{,}5=5{,}5\unite V$. Le voltmètre se connecte aux deux bornes de ce second récepteur.\par
\pcFigureEx{figures/rev-cor-serie-9V.pdf}{Le voltmètre mesure la tension du second récepteur en dérivation ; 3,5 + 5,5 = 9 V.}
\end{corrige}
\end{minipage}
\par\noindent\begin{minipage}{\linewidth}
\begin{corrige}{8}{Tensions en dérivation}
\begin{questions}
\item $U_1=U_2=12\unite V$.
\item Non ; elles dépendent des dipôles.
\end{questions}
\pcFigureEx{figures/rev-cor-resistances-12.pdf}{La tension commune ne force pas les intensités à être égales si les résistances diffèrent.}
\end{corrige}
\end{minipage}
\par\noindent\begin{minipage}{\linewidth}
\begin{corrige}{9}{Une inversion}
\begin{questions}
\item Les bornes V et COM peuvent être inversées par rapport à la convention choisie.
\item Non ; le signe de cette mesure n’impose pas un courant nul.
\end{questions}
\pcFigureEx{figures/cor-voltmetre-inverse.pdf}{Inverser les bornes de mesure inverse le signe affiché, sans imposer un courant nul dans le récepteur.}
\end{corrige}
\end{minipage}
\par\noindent\begin{minipage}{\linewidth}
\begin{corrige}{10}{Commander séparément}
\begin{questions}
\item Deux branches en dérivation.
\item Un interrupteur en série avec sa lampe dans chaque branche.
\end{questions}
\pcFigureEx{figures/cor-commandes-independantes.pdf}{Chaque interrupteur se place en série avec sa lampe dans sa propre branche.}
\end{corrige}
\end{minipage}
\par\noindent\begin{minipage}{\linewidth}
\begin{corrige}{11}{Une mesure incohérente}
\begin{questions}
\item Non ; il faut examiner protocole, stabilité et appareils.
\item Vérifier les connexions et calibres, puis mesurer dans les mêmes conditions avec un circuit stable.
\end{questions}
\end{corrige}
\end{minipage}
\par\noindent\begin{minipage}{\linewidth}
\begin{corrige}{12}{Un nœud à trois sorties}
\begin{questions}
\item $0{,}35\unite A$.
\item Non : elles sont égales dans le modèle du nœud.
\end{questions}
\pcFigureEx{figures/rev-cor-noeud-090.pdf}{0,90 − 0,20 − 0,35 = 0,35 A ; contrôler la somme des trois sorties.}
\end{corrige}
\end{minipage}
\par\noindent\begin{minipage}{\linewidth}
\begin{corrige}{13}{Une lampe retirée}
\begin{questions}
\item Oui, son chemin fermé vers la source demeure.
\item Une source idéale ou suffisamment stable dans les conditions du montage.
\end{questions}
\pcFigureEx{figures/cor-branche-ouverte.pdf}{L’autre branche conserve un trajet fermé ; une tension exactement inchangée suppose la source idéale du modèle.}
\end{corrige}
\end{minipage}
\par\noindent\begin{minipage}{\linewidth}
\begin{corrige}{14}{Court-circuit}
\begin{questions}
\item Le fil peut imposer une tension très faible à la lampe et créer une circulation de courant importante dans le circuit.
\item Couper l’alimentation et faire vérifier le montage.
\end{questions}
\pcFigureEx{figures/cor-lampe-court-circuitee.pdf}{Le fil contourne le récepteur ; couper l’alimentation avant toute modification.}
\end{corrige}
\end{minipage}
\par\noindent\begin{minipage}{\linewidth}
\begin{corrige}{15}{Un cahier des charges}
\begin{questions}
\item Dérivation sur 6 V ; source capable de fournir au moins 0,30 A dans ce modèle.
\item En série, les tensions se partagent et le courant est commun : ces conditions ne donnent pas les deux fonctionnements demandés.
\end{questions}
\pcFigureEx{figures/rev-cor-eclairage-6V.pdf}{Les deux lampes reçoivent chacune 6 V et les intensités de branches s’ajoutent au générateur.}
\end{corrige}
\end{minipage}
\par\noindent\begin{minipage}{\linewidth}
\phantomsection\label{cor-defi-13}\pcSolutionDefi{13}{Le montage à vérifier}\begin{questions}
\item La première branche totalise $4+5=9\unite V$ ; la résistance directement en dérivation reçoit aussi 9 V.
\item $I_G=0{,}150+0{,}060=0{,}210\unite A$.
\item Voltmètre aux bornes de la résistance ; ampèremètre inséré en série dans la première branche, alimentation coupée pendant la modification et contrôle des bornes/calibres.
\end{questions}
\par\smallskip\noindent\textbf{Pour comprendre.}\enspace Contrôler séparément chaque branche : $4+5=9\unite V$ pour les lampes et 9 V pour la résistance. Convertir les courants avant la somme : $0{,}150+0{,}060=0{,}210\unite A$.\par
\pcFigureEx{figures/rev-cor-montage-9V.pdf}{Une branche contient deux lampes en série ; A se place dans cette branche, V aux bornes de la résistance.}
\end{minipage}
\par\noindent\begin{minipage}{\linewidth}
\phantomsection\label{cor-activite-13}\pcSousTitre{Pistes pour l’investigation}
En série, les intensités doivent être compatibles entre elles. En dérivation, IG doit être compatible avec I1 + I2 ; les tensions entre les mêmes nœuds sont égales. Les écarts se discutent avec la stabilité de la source et la précision des appareils. Aucune mesure n’exige de brancher un ampèremètre directement sur la source.\par\end{minipage}
